% Part: counterfactuals
% Chapter: minimal-change-semantics
% Section: transitivity

\documentclass[../../../include/open-logic-section]{subfiles}

\begin{document}

\olfileid[pt]{cnt}{min}{tra}

\olsection{Transitividade}

Para o condicional material, vale a regra da cadeia: $!A \lif !B, !B
\lif !C \Entails !A \lif !C$. Em outras palavras, o condicional material
é transitivo. O mesmo vale para os contrafactuais? Considere o
seguinte exemplo de Stalnaker.
\begin{quote}
  Se J.~Edgar Hoover tivesse nascido russo, teria sido comunista.

  Se J.~Edgar Hoover fosse comunista, teria sido traidor.

  Portanto, se J.~Edgar Hoover tivesse nascido russo, teria
  sido traidor.
\end{quote}
Se Hoover tivesse nascido (na mesma época em que de fato nasceu), não nos
Estados Unidos, mas na Rússia, teria crescido na União Soviética
e se tornado comunista (suponhamos). Portanto, a primeira premissa é
verdadeira. Do mesmo modo, a segunda premissa, considerada isoladamente, é
verdadeira. A conclusão, contudo, é falsa: com toda probabilidade, Hoover
teria sido um comunista fervoroso se tivesse nascido na URSS,
e não um traidor de seu país. A atribuição intuitiva de
valores de verdade é confirmada pela explicação de Stalnaker--Lewis. O mundo
possível mais próximo do nosso em que a única mudança é o local de
nascimento de Hoover é aquele em que Hoover cresce como bom cidadão da
URSS. Esse é o mundo possível mais próximo no qual o antecedente da
primeira premissa e o da conclusão são verdadeiros; nesse mundo, Hoover
é membro leal do Partido Comunista e, portanto, não é traidor. Para
avaliar a segunda premissa, porém, precisamos examinar um mundo diferente:
o mundo mais próximo em que Hoover é comunista, no qual
nasceu nos Estados Unidos, mudou de lado e, assim, tornou-se
traidor.\footnote{É claro que, para compreender a força do exemplo,
  precisamos aceitar algumas suposições metafísicas e políticas;
  p.ex., que seria possível Hoover ter nascido de
  pais russos, ou que comunistas nos Estados Unidos da década de 1950 eram
  traidores de seu país.}

\begin{prob}
  Encontre um exemplo intuitivo e convincente da falha da transitividade
  dos contrafactuais.
\end{prob}

\begin{ex}\ollabel{ex:trans-counterex}
  A semântica de esferas invalida a inferência, O.s.\@, temos $p
  \cif q, q \cif r \Entails/ p \cif r$. Considere o modelo $\mModel{M}
  = \tuple{W, O, V}$, em que $W = \{w, w_1, w_2\}$, $O_w = \{\{w\}, \{w,
  w_1\}, \{w, w_1, w_2\}\}$, $V(p) = \{w_2\}$, $V(q) = \{w_1, w_2\}$,
  e $V(r) = \{w_1\}$. Existe uma esfera que admite $p$, $S = \{w, w_1,
  w_2\}$, e $p \lif q$ é verdadeiro em todos os mundos nela, portanto $\mSat{M}{p
    \cif q}[w]$. Existe também uma esfera que admite $q$, $S' = \{w,
  w_1\}$, e $\mSat/{M}{q \lif r}$ é verdadeiro em todos os mundos nela, portanto
  $\mSat{M}{q \cif r}[w]$. Contudo, a esfera que admite $p$, $\{w,
  w_1, w_2\}$, contém um mundo, a saber,~$w_2$, no qual $\mSat/{M}{p \lif
    r}[w_2]$.
\end{ex}

\begin{prob}
Desenhe o diagrama de esferas correspondente ao contraexemplo em
\olref[cnt][min][tra]{ex:trans-counterex}.
\end{prob}

\begin{prob}
  Em \olref[cnt][min][tra]{ex:trans-counterex}, o mundo $w_2$ é aquele em que
  Hoover nasce na Rússia, é comunista e não é traidor; e
  $w_1$ é o mundo em que Hoover nasce nos Estados Unidos, é comunista
  e é traidor. Nesse modelo, $w_1$ está mais próximo de~$w$ do que
  $w_2$. Isso é necessário? Você pode dar um contraexemplo que
  não suponha que Hoover nascer na Rússia seja uma possibilidade mais remota
  do que ele ser comunista?
\end{prob}

\end{document}
